\scanpage{13}
structure complexe associée à l'orientation (dans le cas où une telle orientation est choisie). Noter que dans les cas 1\textsuperscript{o} et 3\textsuperscript{o}, où les similitudes sont des isométries, on peut \uline{normaliser} la métrique -- on parlera alors de \textquotedblleft surface standard normalisée\textquotedblright\ -- déterminée, pour chacun des \deleted{\unsure{deux|trois}} trois types, à \uline{isom} près.

Une carte cellulaire sur une surface riemannienne \inserted{$X$} est dite \uline{géodésique} si \deleted{elle} ses arêtes ouvertes sont des \inserted{arcs de} géodésiques ; elle est dite \uline{métriquement régulière} si elle est géodésique, et si le groupe des automorphismes \uline{isométriques} de la carte est transitif sur l'ens des répères (donc simplement transitif sur l'ens des répères) On dit que c'est une \uline{carte géodésique \unsure{standart|standard}} si elle est métriquement régulière, \inserted{complète (cf. 1\textsuperscript{o}),} et si \deleted{\unsure{la}} $X$ \deleted{\unsure{y}} est une surface riemannienne \deleted{\unsure{normalisée}} \inserted{standart} \unsure{?} ]

\deleted{On dit qu'elle est} \deleted{\unsure{de si|de type}} \inserted{Ceci posé, on doit avoir ceci} ]

a) Pour toute carte standart \inserted{$(X\supset K\supset S)$}, il existe \deleted{\unsure{sur|une}} une carte géodésique standard $(X_0,K_0,S_0)$ qui lui est (topologiquement) isomorphe.

\deleted{\unsure{Bi}} b) Pour deux cartes géodésiques standard, dans tte classe d'isotopie d'isomorphismes
% archive sheet number: 12

